
@@M0@@. En
itérant l'opération @@M1@@ fois, on arrive à @@M2@@, qui est nul, et on a donc écrit @@M3@@ comme un élément de @@M4@@. Compte tenu de (6.6.5), les deux filtrations sont donc identiques.

6.7. On suppose maintenant que @@M5@@ est un @@M6@@-anneau augmenté vers @@M7@@, l'homomorphisme d'augmentation étant noté @@M8@@. On dispose donc sur @@M9@@ de la filtration canonique de @@M10@@-anneau augmenté, et on peut donc associer à @@M11@@ son anneau de Chern.

Si @@M12@@, alors @@M13@@ par définition de la filtration de @@M14@@-anneau augmenté. On note @@M15@@ et on appelle \uline{classe de Chern totale de} @@M16@@ l'image de @@M17@@ dans @@M18@@~; on note @@M19@@, et on appelle @@M20@@-ième classe de Chern de @@M21@@, le @@M22@@-ième coefficient de @@M23@@~; on a donc :
@@M24@@
On pose alors :
@@M25@@
et on appelle @@M26@@ la classe de Chern complétée de @@M27@@.

\textbf{Proposition 6.8.} \uline{La classe de Chern complétée} @@M28@@, \uline{est} \uline{un} @@M29@@-@@M30@@-\uline{homomorphisme de} @@M31@@-@@M32@@-\uline{algèbres augmentées}.

Il est clair que l'homomorphisme ainsi défini est additif~; pour vérifier qu'il est multiplicatif, on écrit :
@@M33@@
